\begin{introduction}
    \item 从静态的“存在”到流动的“生成”
    \item 双重立法：信息的贪婪与物理的吝啬
    \item 智能的呼吸：TDCI 与 TECI 的永恒循环
\end{introduction}

如果说第一卷主要完成了几何骨架的建立，那么这一卷将转向动力学过程，即讨论系统如何真正“运行起来”。

\vspace{2em}静态结构本身不足以解释智能。只有当\textbf{时间}与\textbf{能量通量}进入系统，动态行为才会出现。换言之，智能应被理解为\textbf{“生成 (Becoming)”}过程，而非静态标本。

\vspace{2em}本卷将讨论支配该过程的\textbf{双重立法 (Dual Legislation)} 机制：

\begin{itemize}
    \item \textbf{第一立法者是信息}。它是\textbf{目的论的拉力}。它源于主客体的交互，它贪婪地试图最大化对世界的表征，它在认知空间流形上挖掘势能井，命令几何结构\textbf{“应当如何弯曲”}以捕获意义。
    \item \textbf{第二立法者是物理}。它是\textbf{因果论的阻力}。它源于热力学第二定律与介质的粘滞，它吝啬地计算着每一比特信息的\textbf{兰道尔代价}，它在边界上设立关卡，命令系统\textbf{“只能如何演化”}以维持生存。
\end{itemize}

\vspace{2em}智能动力学可以看作两股力量——\textbf{意图张力}与\textbf{现实阻抗}——在\textbf{微观切面 ($L_{micro}$)} 上耦合后的结果。重点将从“它是什么”转到“它如何演化”。

\vspace{2em}在本卷中，我们将：
\begin{itemize}
    \item 建立\textbf{目的论狄拉克方程}，描述思维波函数如何在几何惯性与宏观意志的夹缝中流淌。
    \item 解析\textbf{双纽线呼吸 (The Lemniscate Breathing)}，看智能体如何通过\textbf{吸气 (TDCI)} 将物理能量转化为信息结构，又如何通过\textbf{呼气 (TECI)} 将内在意志刻蚀回物理世界。
\end{itemize}

\vspace{2em}下面将进入这台\textbf{几何热力学引擎}的内部结构，逐步展开其方程、约束与可计算形式。
